% =============================================================================
% _tikz-style.tex --- FLUX_CNN 论文 TikZ 共享样式包
%
% 用法（每张 figX-Y.tex 头部 \input 本文件）:
%   \documentclass[border=4pt,varwidth=18cm]{standalone}
%   % =============================================================================
% _tikz-style.tex --- FLUX_CNN 论文 TikZ 共享样式包
%
% 用法（每张 figX-Y.tex 头部 \input 本文件）:
%   \documentclass[border=4pt,varwidth=18cm]{standalone}
%   % =============================================================================
% _tikz-style.tex --- FLUX_CNN 论文 TikZ 共享样式包
%
% 用法（每张 figX-Y.tex 头部 \input 本文件）:
%   \documentclass[border=4pt,varwidth=18cm]{standalone}
%   % =============================================================================
% _tikz-style.tex --- FLUX_CNN 论文 TikZ 共享样式包
%
% 用法（每张 figX-Y.tex 头部 \input 本文件）:
%   \documentclass[border=4pt,varwidth=18cm]{standalone}
%   \input{_tikz-style.tex}
%   \begin{document}
%   \begin{tikzpicture}[<额外局部样式>]
%     ...
%   \end{tikzpicture}
%   \end{document}
%
% 风格基线: IEEE 期刊配色 / 白底黑字 / Times New Roman + Microsoft YaHei /
%           禁卡通 / 禁渐变 / 禁阴影 / 禁 3D / 禁 emoji
% 与 figures-rendered/_style.py (matplotlib) 配色一致, 保证全文图风格统一.
% =============================================================================

% ---- 字体（XeLaTeX 必须）----
% 中文字体：Microsoft YaHei（保持原配置，本次只调整英文）.
% 英文字体：Times New Roman, 把 \sffamily / \rmfamily 同时映射到 Times,
%           保证 TikZ 节点无论用哪个 family, 拉丁字符都是 Times New Roman.
% 数学公式（mathmode）保留 latin modern math, 与 Times 视觉差异可接受.
\usepackage{amsmath}
\usepackage{amssymb}
\usepackage{xeCJK}
\setCJKmainfont{Microsoft YaHei}
\setCJKsansfont{Microsoft YaHei}
\setmainfont{Times New Roman}
\setsansfont{Times New Roman}

% ---- 颜色（与 figures-rendered/_style.py PALETTE='default' 同款）----
\usepackage{xcolor}
\definecolor{coBaseline}{HTML}{1f4e79}  % 深蓝 - 主体强调 / 基线
\definecolor{coImproved}{HTML}{2e7d32}  % 深绿 - 提升 / 高亮
\definecolor{coThird}   {HTML}{b8860b}  % 深金 - 第三系列
\definecolor{coFourth}  {HTML}{6a4c93}  % 深紫 - 第四系列
\definecolor{coAnno}    {HTML}{404040}  % 学术深灰 - 注记
\definecolor{coGuide}   {HTML}{808080}  % 辅助线
\definecolor{coGrid}    {HTML}{cccccc}  % 网格
\definecolor{coAccent}  {HTML}{c00000}  % 强调红 - 仅用于关键警示 / 跨节点
% 浅色填充（CMYK 转 RGB 近似）
\definecolor{flBlue}    {RGB}{226,238,247}  % 淡蓝 CMYK 30/10/0/0 ≈
\definecolor{flYellow}  {RGB}{252,242,219}  % 淡黄 CMYK 0/10/30/0 ≈
\definecolor{flGray}    {RGB}{235,235,235}  % 淡灰 CMYK 0/0/0/10 ≈
\definecolor{flGreen}   {RGB}{226,239,224}  % 浅绿 - 高亮区域

% ---- TikZ + 常用库 ----
\usepackage{tikz}
\usetikzlibrary{
    arrows.meta,
    positioning,
    calc,
    fit,
    backgrounds,
    decorations.pathreplacing,
    decorations.pathmorphing,
    shapes.geometric,
    shapes.misc,
    shapes.multipart,
    matrix,
    chains,
    shadows.blur,    % 不用阴影 (禁用), 但 fit 节点偶尔用 inner sep
    patterns
}

% ---- 全局 TikZ 默认 ----
\tikzset{
    every picture/.style={
        line cap=round,
        line join=round,
        >={Stealth[length=2.4mm, width=1.6mm]},
        font=\sffamily\small,
    },
    % ---- 节点基础类型 ----
    %
    % 主体方框 (深色边框, 白底) - 用于一般模块
    main box/.style={
        draw=coAnno,
        line width=0.7pt,
        rectangle,
        rounded corners=1.5pt,
        minimum width=20mm,
        minimum height=8mm,
        align=center,
        inner sep=2pt,
    },
    % 强调方框 (深蓝填充 alpha)
    accent box/.style={
        main box,
        fill=flBlue,
    },
    % 计算单元 / PE (淡黄填充)
    pe box/.style={
        main box,
        fill=flYellow,
    },
    % SRAM / 存储 (淡灰填充)
    mem box/.style={
        main box,
        fill=flGray,
    },
    % 控制单元 / FSM (浅绿填充)
    ctrl box/.style={
        main box,
        fill=flGreen,
    },
    % 大容器 (圆角矩形, 虚线边框, 用于 multicore_top / NUMA 节点等)
    container box/.style={
        draw=coAnno,
        line width=0.7pt,
        rounded corners=2pt,
        rectangle,
        dashed,
        inner sep=3mm,
    },
    container solid/.style={
        container box,
        solid,
    },
    %
    % ---- 箭头类型 ----
    %
    flow/.style={
        -{Stealth[length=2.4mm,width=1.6mm]},
        line width=0.7pt,
        draw=coAnno,
    },
    flow thick/.style={
        flow,
        line width=1.2pt,
    },
    flow dashed/.style={
        flow,
        dashed,
    },
    flow accent/.style={
        -{Stealth[length=2.4mm,width=1.6mm]},
        line width=1.0pt,
        draw=coAccent,
        dashed,
    },
    flow data/.style={
        -{Stealth[length=2.4mm,width=1.6mm]},
        line width=0.9pt,
        draw=coBaseline,
    },
    flow weight/.style={
        -{Stealth[length=2.4mm,width=1.6mm]},
        line width=0.9pt,
        draw=coImproved,
    },
    flow psum/.style={
        -{Stealth[length=2.4mm,width=1.6mm]},
        line width=0.9pt,
        draw=coThird,
    },
    bus/.style={
        -{Stealth[length=3mm,width=2mm]},
        line width=1.6pt,
        draw=coAnno,
    },
    handshake/.style={
        <->,
        line width=0.7pt,
        draw=coAnno,
    },
    %
    % ---- 标注 / 文字 ----
    %
    label small/.style={
        font=\sffamily\footnotesize,
        text=coAnno,
        align=center,
    },
    label tiny/.style={
        font=\sffamily\scriptsize,
        text=coAnno,
        align=center,
    },
    label bold/.style={
        font=\sffamily\bfseries\small,
        text=black,
        align=center,
    },
    formula/.style={
        font=\small,
        align=center,
    },
    % 边线上的标注 (背景白底, 不挡线)
    edge label/.style={
        font=\sffamily\scriptsize,
        text=coAnno,
        fill=white,
        inner sep=1.2pt,
        align=center,
    },
    edge label accent/.style={
        edge label,
        text=coAccent,
        font=\sffamily\bfseries\scriptsize,
    },
    %
    % ---- 图例 ----
    %
    legend box/.style={
        draw=coAnno,
        line width=0.4pt,
        fill=white,
        rounded corners=1pt,
        inner sep=2.5pt,
        font=\sffamily\scriptsize,
        text=coAnno,
        align=left,
    },
    %
    % ---- 网格 / 阵列 (用于 PE 阵列示意) ----
    %
    grid cell/.style={
        draw=coAnno,
        line width=0.3pt,
        minimum width=4mm,
        minimum height=4mm,
        inner sep=0pt,
        rectangle,
    },
    grid cell light/.style={
        grid cell,
        fill=flGray,
    },
    grid cell accent/.style={
        grid cell,
        fill=flBlue,
    },
    %
    % ---- 时序波形 (用于内嵌时序示意时, 非主要; 主要用 WaveDrom) ----
    %
    wave hi/.style={line width=1pt, draw=coBaseline},
    wave lo/.style={line width=1pt, draw=coBaseline},
}

% ---- 通用辅助宏 ----

% 章节图标题命令 (备选, 一般由论文 caption 承担, 不在 TikZ 内画)
% \newcommand{\figtitle}[1]{\node[font=\bfseries\small,text=black] {#1};}

% 端口标记: 在节点边缘画个小圆点表示端口
\newcommand{\port}[2][black]{\fill[#1] (#2) circle (0.5mm);}

% 中文/英文双行标签: \biLabel{中文}{English}
\newcommand{\biLabel}[2]{\shortstack{\sffamily #1\\\itshape\footnotesize #2}}

% =============================================================================
% 使用约定 (写 figX-Y.tex 时遵守):
%
% 1. 不在 TikZ 内画"图 X.Y"标题——由 Word caption 承担
% 2. 节点用上述预定义样式 (main box / accent box / pe box / mem box / ctrl box / container box)
%    避免每张图重复定义节点样式
% 3. 数据流统一用 flow data (蓝) / flow weight (绿) / flow psum (金) 三色编码
%    握手 / 控制流用 flow / flow dashed (灰)
%    强调 / 警示用 flow accent (红虚线), 慎用
% 4. 文字标注用 label small / label tiny / label bold / edge label
% 5. 图例统一放右下角, 用 legend box 样式
% 6. 中英混排时: 中文用 \sffamily (思源黑体), 英文数学用 \rmfamily (Times New Roman)
% =============================================================================

%   \begin{document}
%   \begin{tikzpicture}[<额外局部样式>]
%     ...
%   \end{tikzpicture}
%   \end{document}
%
% 风格基线: IEEE 期刊配色 / 白底黑字 / Times New Roman + Microsoft YaHei /
%           禁卡通 / 禁渐变 / 禁阴影 / 禁 3D / 禁 emoji
% 与 figures-rendered/_style.py (matplotlib) 配色一致, 保证全文图风格统一.
% =============================================================================

% ---- 字体（XeLaTeX 必须）----
% 中文字体：Microsoft YaHei（保持原配置，本次只调整英文）.
% 英文字体：Times New Roman, 把 \sffamily / \rmfamily 同时映射到 Times,
%           保证 TikZ 节点无论用哪个 family, 拉丁字符都是 Times New Roman.
% 数学公式（mathmode）保留 latin modern math, 与 Times 视觉差异可接受.
\usepackage{amsmath}
\usepackage{amssymb}
\usepackage{xeCJK}
\setCJKmainfont{Microsoft YaHei}
\setCJKsansfont{Microsoft YaHei}
\setmainfont{Times New Roman}
\setsansfont{Times New Roman}

% ---- 颜色（与 figures-rendered/_style.py PALETTE='default' 同款）----
\usepackage{xcolor}
\definecolor{coBaseline}{HTML}{1f4e79}  % 深蓝 - 主体强调 / 基线
\definecolor{coImproved}{HTML}{2e7d32}  % 深绿 - 提升 / 高亮
\definecolor{coThird}   {HTML}{b8860b}  % 深金 - 第三系列
\definecolor{coFourth}  {HTML}{6a4c93}  % 深紫 - 第四系列
\definecolor{coAnno}    {HTML}{404040}  % 学术深灰 - 注记
\definecolor{coGuide}   {HTML}{808080}  % 辅助线
\definecolor{coGrid}    {HTML}{cccccc}  % 网格
\definecolor{coAccent}  {HTML}{c00000}  % 强调红 - 仅用于关键警示 / 跨节点
% 浅色填充（CMYK 转 RGB 近似）
\definecolor{flBlue}    {RGB}{226,238,247}  % 淡蓝 CMYK 30/10/0/0 ≈
\definecolor{flYellow}  {RGB}{252,242,219}  % 淡黄 CMYK 0/10/30/0 ≈
\definecolor{flGray}    {RGB}{235,235,235}  % 淡灰 CMYK 0/0/0/10 ≈
\definecolor{flGreen}   {RGB}{226,239,224}  % 浅绿 - 高亮区域

% ---- TikZ + 常用库 ----
\usepackage{tikz}
\usetikzlibrary{
    arrows.meta,
    positioning,
    calc,
    fit,
    backgrounds,
    decorations.pathreplacing,
    decorations.pathmorphing,
    shapes.geometric,
    shapes.misc,
    shapes.multipart,
    matrix,
    chains,
    shadows.blur,    % 不用阴影 (禁用), 但 fit 节点偶尔用 inner sep
    patterns
}

% ---- 全局 TikZ 默认 ----
\tikzset{
    every picture/.style={
        line cap=round,
        line join=round,
        >={Stealth[length=2.4mm, width=1.6mm]},
        font=\sffamily\small,
    },
    % ---- 节点基础类型 ----
    %
    % 主体方框 (深色边框, 白底) - 用于一般模块
    main box/.style={
        draw=coAnno,
        line width=0.7pt,
        rectangle,
        rounded corners=1.5pt,
        minimum width=20mm,
        minimum height=8mm,
        align=center,
        inner sep=2pt,
    },
    % 强调方框 (深蓝填充 alpha)
    accent box/.style={
        main box,
        fill=flBlue,
    },
    % 计算单元 / PE (淡黄填充)
    pe box/.style={
        main box,
        fill=flYellow,
    },
    % SRAM / 存储 (淡灰填充)
    mem box/.style={
        main box,
        fill=flGray,
    },
    % 控制单元 / FSM (浅绿填充)
    ctrl box/.style={
        main box,
        fill=flGreen,
    },
    % 大容器 (圆角矩形, 虚线边框, 用于 multicore_top / NUMA 节点等)
    container box/.style={
        draw=coAnno,
        line width=0.7pt,
        rounded corners=2pt,
        rectangle,
        dashed,
        inner sep=3mm,
    },
    container solid/.style={
        container box,
        solid,
    },
    %
    % ---- 箭头类型 ----
    %
    flow/.style={
        -{Stealth[length=2.4mm,width=1.6mm]},
        line width=0.7pt,
        draw=coAnno,
    },
    flow thick/.style={
        flow,
        line width=1.2pt,
    },
    flow dashed/.style={
        flow,
        dashed,
    },
    flow accent/.style={
        -{Stealth[length=2.4mm,width=1.6mm]},
        line width=1.0pt,
        draw=coAccent,
        dashed,
    },
    flow data/.style={
        -{Stealth[length=2.4mm,width=1.6mm]},
        line width=0.9pt,
        draw=coBaseline,
    },
    flow weight/.style={
        -{Stealth[length=2.4mm,width=1.6mm]},
        line width=0.9pt,
        draw=coImproved,
    },
    flow psum/.style={
        -{Stealth[length=2.4mm,width=1.6mm]},
        line width=0.9pt,
        draw=coThird,
    },
    bus/.style={
        -{Stealth[length=3mm,width=2mm]},
        line width=1.6pt,
        draw=coAnno,
    },
    handshake/.style={
        <->,
        line width=0.7pt,
        draw=coAnno,
    },
    %
    % ---- 标注 / 文字 ----
    %
    label small/.style={
        font=\sffamily\footnotesize,
        text=coAnno,
        align=center,
    },
    label tiny/.style={
        font=\sffamily\scriptsize,
        text=coAnno,
        align=center,
    },
    label bold/.style={
        font=\sffamily\bfseries\small,
        text=black,
        align=center,
    },
    formula/.style={
        font=\small,
        align=center,
    },
    % 边线上的标注 (背景白底, 不挡线)
    edge label/.style={
        font=\sffamily\scriptsize,
        text=coAnno,
        fill=white,
        inner sep=1.2pt,
        align=center,
    },
    edge label accent/.style={
        edge label,
        text=coAccent,
        font=\sffamily\bfseries\scriptsize,
    },
    %
    % ---- 图例 ----
    %
    legend box/.style={
        draw=coAnno,
        line width=0.4pt,
        fill=white,
        rounded corners=1pt,
        inner sep=2.5pt,
        font=\sffamily\scriptsize,
        text=coAnno,
        align=left,
    },
    %
    % ---- 网格 / 阵列 (用于 PE 阵列示意) ----
    %
    grid cell/.style={
        draw=coAnno,
        line width=0.3pt,
        minimum width=4mm,
        minimum height=4mm,
        inner sep=0pt,
        rectangle,
    },
    grid cell light/.style={
        grid cell,
        fill=flGray,
    },
    grid cell accent/.style={
        grid cell,
        fill=flBlue,
    },
    %
    % ---- 时序波形 (用于内嵌时序示意时, 非主要; 主要用 WaveDrom) ----
    %
    wave hi/.style={line width=1pt, draw=coBaseline},
    wave lo/.style={line width=1pt, draw=coBaseline},
}

% ---- 通用辅助宏 ----

% 章节图标题命令 (备选, 一般由论文 caption 承担, 不在 TikZ 内画)
% \newcommand{\figtitle}[1]{\node[font=\bfseries\small,text=black] {#1};}

% 端口标记: 在节点边缘画个小圆点表示端口
\newcommand{\port}[2][black]{\fill[#1] (#2) circle (0.5mm);}

% 中文/英文双行标签: \biLabel{中文}{English}
\newcommand{\biLabel}[2]{\shortstack{\sffamily #1\\\itshape\footnotesize #2}}

% =============================================================================
% 使用约定 (写 figX-Y.tex 时遵守):
%
% 1. 不在 TikZ 内画"图 X.Y"标题——由 Word caption 承担
% 2. 节点用上述预定义样式 (main box / accent box / pe box / mem box / ctrl box / container box)
%    避免每张图重复定义节点样式
% 3. 数据流统一用 flow data (蓝) / flow weight (绿) / flow psum (金) 三色编码
%    握手 / 控制流用 flow / flow dashed (灰)
%    强调 / 警示用 flow accent (红虚线), 慎用
% 4. 文字标注用 label small / label tiny / label bold / edge label
% 5. 图例统一放右下角, 用 legend box 样式
% 6. 中英混排时: 中文用 \sffamily (思源黑体), 英文数学用 \rmfamily (Times New Roman)
% =============================================================================

%   \begin{document}
%   \begin{tikzpicture}[<额外局部样式>]
%     ...
%   \end{tikzpicture}
%   \end{document}
%
% 风格基线: IEEE 期刊配色 / 白底黑字 / Times New Roman + Microsoft YaHei /
%           禁卡通 / 禁渐变 / 禁阴影 / 禁 3D / 禁 emoji
% 与 figures-rendered/_style.py (matplotlib) 配色一致, 保证全文图风格统一.
% =============================================================================

% ---- 字体（XeLaTeX 必须）----
% 中文字体：Microsoft YaHei（保持原配置，本次只调整英文）.
% 英文字体：Times New Roman, 把 \sffamily / \rmfamily 同时映射到 Times,
%           保证 TikZ 节点无论用哪个 family, 拉丁字符都是 Times New Roman.
% 数学公式（mathmode）保留 latin modern math, 与 Times 视觉差异可接受.
\usepackage{amsmath}
\usepackage{amssymb}
\usepackage{xeCJK}
\setCJKmainfont{Microsoft YaHei}
\setCJKsansfont{Microsoft YaHei}
\setmainfont{Times New Roman}
\setsansfont{Times New Roman}

% ---- 颜色（与 figures-rendered/_style.py PALETTE='default' 同款）----
\usepackage{xcolor}
\definecolor{coBaseline}{HTML}{1f4e79}  % 深蓝 - 主体强调 / 基线
\definecolor{coImproved}{HTML}{2e7d32}  % 深绿 - 提升 / 高亮
\definecolor{coThird}   {HTML}{b8860b}  % 深金 - 第三系列
\definecolor{coFourth}  {HTML}{6a4c93}  % 深紫 - 第四系列
\definecolor{coAnno}    {HTML}{404040}  % 学术深灰 - 注记
\definecolor{coGuide}   {HTML}{808080}  % 辅助线
\definecolor{coGrid}    {HTML}{cccccc}  % 网格
\definecolor{coAccent}  {HTML}{c00000}  % 强调红 - 仅用于关键警示 / 跨节点
% 浅色填充（CMYK 转 RGB 近似）
\definecolor{flBlue}    {RGB}{226,238,247}  % 淡蓝 CMYK 30/10/0/0 ≈
\definecolor{flYellow}  {RGB}{252,242,219}  % 淡黄 CMYK 0/10/30/0 ≈
\definecolor{flGray}    {RGB}{235,235,235}  % 淡灰 CMYK 0/0/0/10 ≈
\definecolor{flGreen}   {RGB}{226,239,224}  % 浅绿 - 高亮区域

% ---- TikZ + 常用库 ----
\usepackage{tikz}
\usetikzlibrary{
    arrows.meta,
    positioning,
    calc,
    fit,
    backgrounds,
    decorations.pathreplacing,
    decorations.pathmorphing,
    shapes.geometric,
    shapes.misc,
    shapes.multipart,
    matrix,
    chains,
    shadows.blur,    % 不用阴影 (禁用), 但 fit 节点偶尔用 inner sep
    patterns
}

% ---- 全局 TikZ 默认 ----
\tikzset{
    every picture/.style={
        line cap=round,
        line join=round,
        >={Stealth[length=2.4mm, width=1.6mm]},
        font=\sffamily\small,
    },
    % ---- 节点基础类型 ----
    %
    % 主体方框 (深色边框, 白底) - 用于一般模块
    main box/.style={
        draw=coAnno,
        line width=0.7pt,
        rectangle,
        rounded corners=1.5pt,
        minimum width=20mm,
        minimum height=8mm,
        align=center,
        inner sep=2pt,
    },
    % 强调方框 (深蓝填充 alpha)
    accent box/.style={
        main box,
        fill=flBlue,
    },
    % 计算单元 / PE (淡黄填充)
    pe box/.style={
        main box,
        fill=flYellow,
    },
    % SRAM / 存储 (淡灰填充)
    mem box/.style={
        main box,
        fill=flGray,
    },
    % 控制单元 / FSM (浅绿填充)
    ctrl box/.style={
        main box,
        fill=flGreen,
    },
    % 大容器 (圆角矩形, 虚线边框, 用于 multicore_top / NUMA 节点等)
    container box/.style={
        draw=coAnno,
        line width=0.7pt,
        rounded corners=2pt,
        rectangle,
        dashed,
        inner sep=3mm,
    },
    container solid/.style={
        container box,
        solid,
    },
    %
    % ---- 箭头类型 ----
    %
    flow/.style={
        -{Stealth[length=2.4mm,width=1.6mm]},
        line width=0.7pt,
        draw=coAnno,
    },
    flow thick/.style={
        flow,
        line width=1.2pt,
    },
    flow dashed/.style={
        flow,
        dashed,
    },
    flow accent/.style={
        -{Stealth[length=2.4mm,width=1.6mm]},
        line width=1.0pt,
        draw=coAccent,
        dashed,
    },
    flow data/.style={
        -{Stealth[length=2.4mm,width=1.6mm]},
        line width=0.9pt,
        draw=coBaseline,
    },
    flow weight/.style={
        -{Stealth[length=2.4mm,width=1.6mm]},
        line width=0.9pt,
        draw=coImproved,
    },
    flow psum/.style={
        -{Stealth[length=2.4mm,width=1.6mm]},
        line width=0.9pt,
        draw=coThird,
    },
    bus/.style={
        -{Stealth[length=3mm,width=2mm]},
        line width=1.6pt,
        draw=coAnno,
    },
    handshake/.style={
        <->,
        line width=0.7pt,
        draw=coAnno,
    },
    %
    % ---- 标注 / 文字 ----
    %
    label small/.style={
        font=\sffamily\footnotesize,
        text=coAnno,
        align=center,
    },
    label tiny/.style={
        font=\sffamily\scriptsize,
        text=coAnno,
        align=center,
    },
    label bold/.style={
        font=\sffamily\bfseries\small,
        text=black,
        align=center,
    },
    formula/.style={
        font=\small,
        align=center,
    },
    % 边线上的标注 (背景白底, 不挡线)
    edge label/.style={
        font=\sffamily\scriptsize,
        text=coAnno,
        fill=white,
        inner sep=1.2pt,
        align=center,
    },
    edge label accent/.style={
        edge label,
        text=coAccent,
        font=\sffamily\bfseries\scriptsize,
    },
    %
    % ---- 图例 ----
    %
    legend box/.style={
        draw=coAnno,
        line width=0.4pt,
        fill=white,
        rounded corners=1pt,
        inner sep=2.5pt,
        font=\sffamily\scriptsize,
        text=coAnno,
        align=left,
    },
    %
    % ---- 网格 / 阵列 (用于 PE 阵列示意) ----
    %
    grid cell/.style={
        draw=coAnno,
        line width=0.3pt,
        minimum width=4mm,
        minimum height=4mm,
        inner sep=0pt,
        rectangle,
    },
    grid cell light/.style={
        grid cell,
        fill=flGray,
    },
    grid cell accent/.style={
        grid cell,
        fill=flBlue,
    },
    %
    % ---- 时序波形 (用于内嵌时序示意时, 非主要; 主要用 WaveDrom) ----
    %
    wave hi/.style={line width=1pt, draw=coBaseline},
    wave lo/.style={line width=1pt, draw=coBaseline},
}

% ---- 通用辅助宏 ----

% 章节图标题命令 (备选, 一般由论文 caption 承担, 不在 TikZ 内画)
% \newcommand{\figtitle}[1]{\node[font=\bfseries\small,text=black] {#1};}

% 端口标记: 在节点边缘画个小圆点表示端口
\newcommand{\port}[2][black]{\fill[#1] (#2) circle (0.5mm);}

% 中文/英文双行标签: \biLabel{中文}{English}
\newcommand{\biLabel}[2]{\shortstack{\sffamily #1\\\itshape\footnotesize #2}}

% =============================================================================
% 使用约定 (写 figX-Y.tex 时遵守):
%
% 1. 不在 TikZ 内画"图 X.Y"标题——由 Word caption 承担
% 2. 节点用上述预定义样式 (main box / accent box / pe box / mem box / ctrl box / container box)
%    避免每张图重复定义节点样式
% 3. 数据流统一用 flow data (蓝) / flow weight (绿) / flow psum (金) 三色编码
%    握手 / 控制流用 flow / flow dashed (灰)
%    强调 / 警示用 flow accent (红虚线), 慎用
% 4. 文字标注用 label small / label tiny / label bold / edge label
% 5. 图例统一放右下角, 用 legend box 样式
% 6. 中英混排时: 中文用 \sffamily (思源黑体), 英文数学用 \rmfamily (Times New Roman)
% =============================================================================

%   \begin{document}
%   \begin{tikzpicture}[<额外局部样式>]
%     ...
%   \end{tikzpicture}
%   \end{document}
%
% 风格基线: IEEE 期刊配色 / 白底黑字 / Times New Roman + Microsoft YaHei /
%           禁卡通 / 禁渐变 / 禁阴影 / 禁 3D / 禁 emoji
% 与 figures-rendered/_style.py (matplotlib) 配色一致, 保证全文图风格统一.
% =============================================================================

% ---- 字体（XeLaTeX 必须）----
% 中文字体：Microsoft YaHei（保持原配置，本次只调整英文）.
% 英文字体：Times New Roman, 把 \sffamily / \rmfamily 同时映射到 Times,
%           保证 TikZ 节点无论用哪个 family, 拉丁字符都是 Times New Roman.
% 数学公式（mathmode）保留 latin modern math, 与 Times 视觉差异可接受.
\usepackage{amsmath}
\usepackage{amssymb}
\usepackage{xeCJK}
\setCJKmainfont{Microsoft YaHei}
\setCJKsansfont{Microsoft YaHei}
\setmainfont{Times New Roman}
\setsansfont{Times New Roman}

% ---- 颜色（与 figures-rendered/_style.py PALETTE='default' 同款）----
\usepackage{xcolor}
\definecolor{coBaseline}{HTML}{1f4e79}  % 深蓝 - 主体强调 / 基线
\definecolor{coImproved}{HTML}{2e7d32}  % 深绿 - 提升 / 高亮
\definecolor{coThird}   {HTML}{b8860b}  % 深金 - 第三系列
\definecolor{coFourth}  {HTML}{6a4c93}  % 深紫 - 第四系列
\definecolor{coAnno}    {HTML}{404040}  % 学术深灰 - 注记
\definecolor{coGuide}   {HTML}{808080}  % 辅助线
\definecolor{coGrid}    {HTML}{cccccc}  % 网格
\definecolor{coAccent}  {HTML}{c00000}  % 强调红 - 仅用于关键警示 / 跨节点
% 浅色填充（CMYK 转 RGB 近似）
\definecolor{flBlue}    {RGB}{226,238,247}  % 淡蓝 CMYK 30/10/0/0 ≈
\definecolor{flYellow}  {RGB}{252,242,219}  % 淡黄 CMYK 0/10/30/0 ≈
\definecolor{flGray}    {RGB}{235,235,235}  % 淡灰 CMYK 0/0/0/10 ≈
\definecolor{flGreen}   {RGB}{226,239,224}  % 浅绿 - 高亮区域

% ---- TikZ + 常用库 ----
\usepackage{tikz}
\usetikzlibrary{
    arrows.meta,
    positioning,
    calc,
    fit,
    backgrounds,
    decorations.pathreplacing,
    decorations.pathmorphing,
    shapes.geometric,
    shapes.misc,
    shapes.multipart,
    matrix,
    chains,
    shadows.blur,    % 不用阴影 (禁用), 但 fit 节点偶尔用 inner sep
    patterns
}

% ---- 全局 TikZ 默认 ----
\tikzset{
    every picture/.style={
        line cap=round,
        line join=round,
        >={Stealth[length=2.4mm, width=1.6mm]},
        font=\sffamily\small,
    },
    % ---- 节点基础类型 ----
    %
    % 主体方框 (深色边框, 白底) - 用于一般模块
    main box/.style={
        draw=coAnno,
        line width=0.7pt,
        rectangle,
        rounded corners=1.5pt,
        minimum width=20mm,
        minimum height=8mm,
        align=center,
        inner sep=2pt,
    },
    % 强调方框 (深蓝填充 alpha)
    accent box/.style={
        main box,
        fill=flBlue,
    },
    % 计算单元 / PE (淡黄填充)
    pe box/.style={
        main box,
        fill=flYellow,
    },
    % SRAM / 存储 (淡灰填充)
    mem box/.style={
        main box,
        fill=flGray,
    },
    % 控制单元 / FSM (浅绿填充)
    ctrl box/.style={
        main box,
        fill=flGreen,
    },
    % 大容器 (圆角矩形, 虚线边框, 用于 multicore_top / NUMA 节点等)
    container box/.style={
        draw=coAnno,
        line width=0.7pt,
        rounded corners=2pt,
        rectangle,
        dashed,
        inner sep=3mm,
    },
    container solid/.style={
        container box,
        solid,
    },
    %
    % ---- 箭头类型 ----
    %
    flow/.style={
        -{Stealth[length=2.4mm,width=1.6mm]},
        line width=0.7pt,
        draw=coAnno,
    },
    flow thick/.style={
        flow,
        line width=1.2pt,
    },
    flow dashed/.style={
        flow,
        dashed,
    },
    flow accent/.style={
        -{Stealth[length=2.4mm,width=1.6mm]},
        line width=1.0pt,
        draw=coAccent,
        dashed,
    },
    flow data/.style={
        -{Stealth[length=2.4mm,width=1.6mm]},
        line width=0.9pt,
        draw=coBaseline,
    },
    flow weight/.style={
        -{Stealth[length=2.4mm,width=1.6mm]},
        line width=0.9pt,
        draw=coImproved,
    },
    flow psum/.style={
        -{Stealth[length=2.4mm,width=1.6mm]},
        line width=0.9pt,
        draw=coThird,
    },
    bus/.style={
        -{Stealth[length=3mm,width=2mm]},
        line width=1.6pt,
        draw=coAnno,
    },
    handshake/.style={
        <->,
        line width=0.7pt,
        draw=coAnno,
    },
    %
    % ---- 标注 / 文字 ----
    %
    label small/.style={
        font=\sffamily\footnotesize,
        text=coAnno,
        align=center,
    },
    label tiny/.style={
        font=\sffamily\scriptsize,
        text=coAnno,
        align=center,
    },
    label bold/.style={
        font=\sffamily\bfseries\small,
        text=black,
        align=center,
    },
    formula/.style={
        font=\small,
        align=center,
    },
    % 边线上的标注 (背景白底, 不挡线)
    edge label/.style={
        font=\sffamily\scriptsize,
        text=coAnno,
        fill=white,
        inner sep=1.2pt,
        align=center,
    },
    edge label accent/.style={
        edge label,
        text=coAccent,
        font=\sffamily\bfseries\scriptsize,
    },
    %
    % ---- 图例 ----
    %
    legend box/.style={
        draw=coAnno,
        line width=0.4pt,
        fill=white,
        rounded corners=1pt,
        inner sep=2.5pt,
        font=\sffamily\scriptsize,
        text=coAnno,
        align=left,
    },
    %
    % ---- 网格 / 阵列 (用于 PE 阵列示意) ----
    %
    grid cell/.style={
        draw=coAnno,
        line width=0.3pt,
        minimum width=4mm,
        minimum height=4mm,
        inner sep=0pt,
        rectangle,
    },
    grid cell light/.style={
        grid cell,
        fill=flGray,
    },
    grid cell accent/.style={
        grid cell,
        fill=flBlue,
    },
    %
    % ---- 时序波形 (用于内嵌时序示意时, 非主要; 主要用 WaveDrom) ----
    %
    wave hi/.style={line width=1pt, draw=coBaseline},
    wave lo/.style={line width=1pt, draw=coBaseline},
}

% ---- 通用辅助宏 ----

% 章节图标题命令 (备选, 一般由论文 caption 承担, 不在 TikZ 内画)
% \newcommand{\figtitle}[1]{\node[font=\bfseries\small,text=black] {#1};}

% 端口标记: 在节点边缘画个小圆点表示端口
\newcommand{\port}[2][black]{\fill[#1] (#2) circle (0.5mm);}

% 中文/英文双行标签: \biLabel{中文}{English}
\newcommand{\biLabel}[2]{\shortstack{\sffamily #1\\\itshape\footnotesize #2}}

% =============================================================================
% 使用约定 (写 figX-Y.tex 时遵守):
%
% 1. 不在 TikZ 内画"图 X.Y"标题——由 Word caption 承担
% 2. 节点用上述预定义样式 (main box / accent box / pe box / mem box / ctrl box / container box)
%    避免每张图重复定义节点样式
% 3. 数据流统一用 flow data (蓝) / flow weight (绿) / flow psum (金) 三色编码
%    握手 / 控制流用 flow / flow dashed (灰)
%    强调 / 警示用 flow accent (红虚线), 慎用
% 4. 文字标注用 label small / label tiny / label bold / edge label
% 5. 图例统一放右下角, 用 legend box 样式
% 6. 中英混排时: 中文用 \sffamily (思源黑体), 英文数学用 \rmfamily (Times New Roman)
% =============================================================================
